\chapter{Uvod}\label{cha:uvod}
Uvodno poglavje naj vsebuje splošni opis oziroma razlago obravnavane tematike. Predstavite izhodišča doktorskega dela, problematiko in njegov pomen. V uvodu ne predstavljajte rezultatov in sklepov.

V tej predlogi je \textbf{v nevsebinskem delu} (do strani xvi) potrebno izpolniti vse informacije, da bodo ustrezale podatkom kandidata oz.~kandidatke in doktorskega dela. Kazalo vsebine se ob prevajanju osveži samodejno. Pri tem upoštevajte navodila glede uporabe ustreznih slogov, ki so navedena v \textbf{prilogi \ref{cha:priloga}}.

V tej predlogi so \textbf{v vsebinskem delu} navodila in primeri za oblikovanje doktorskega dela, tehnične podrobnosti pa so dodatno opisane v \textbf{prilogi \ref{cha:priloga}}. Celotno besedilo je potrebno nadomestiti z besedilom, ki vsebinsko ustreza doktorskemu delu.

V primeru, da je obravnavana tematika tesno povezana s širšim znanstvenim področjem oz.~poznavanjem določenih pojavov, procesov, postopkov, preračunov itn., ki jih doktorsko delo sicer ne obravnava neposredno, so pa pomembni za njegovo razumevanje, se lahko v tem poglavju predstavi omenjeno širšo tematiko.

\section{Struktura doktorskega dela}\label{sec:struktura_dela}
V tem delu se osredotočite na to, kako je delo strukturirano (po poglavjih), torej npr.~kje je predstavljen nek vsebinski sklop in kako so posamezni vsebinski sklopi med seboj povezani.
